\subsection{Service-Oriented Architecture (SOA)}\label{sec:modeling--service-oriented-architecture-soa}

In the previous section, we discussed \textbf{client-server and multi-tier architectures}, where components are organized according to their roles and responsibilities. Now we consider a \textbf{more complex application}, such as an e-commerce platform. It might need to perform several operations:
\begin{multicols}{2}
    \begin{itemize}
        \item Authenticate customers.
        \item Check product availability.
        \item Process payments.
        \item Arrange shipments.
        \item Send notifications.
    \end{itemize}
\end{multicols}
\noindent
One possibility is to implement all these functionalities within a \emph{single application server}. But \textbf{\emph{what if we want each functionality to be developed, deployed, maintained, and potentially reused independently?}} This is the motivation behind \textbf{Service-Oriented Architecture (SOA)}.

\highspace
\begin{flushleft}
    \textcolor{Green3}{\faIcon{puzzle-piece} \textbf{What is Service-Oriented Architecture?}}
\end{flushleft}
\definition{Service-Oriented Architecture (SOA)} is an architectural style\footnote{%
    An architectural style (\autopageref{def:architectural-style}) is a reusable way of organizing distributed components and defining how they interact.
} in which a \hl{distributed application is organized as a \textbf{collection of loosely coupled services} that \textbf{communicate through well-defined interfaces}.}

\highspace
A \definition{Service} is a \hl{self-contained unit of functionality that exposes operations to other components without requiring them to understand its internal implementation.}

\highspace
\begin{examplebox}[: Payment Service]
    For example, a payment service might expose:
\begin{lstlisting}
PaymentService
    - authorizePayment(amount, currency, cardInfo)
    - capturePayment(transactionId)
    - refundPayment(transactionId)
\end{lstlisting}
    Other applications can \textbf{invoke} these operations \textbf{without knowing}:
    \begin{itemize}
        \item Which programming language implements the service.
        \item Which operating system hosts it.
        \item How its internal algorithms work.
        \item How it stores its data.
    \end{itemize}
\end{examplebox}

\highspace
In general, external applications only need to understand the service's \textbf{public contract}, which specifies the available operations and how to invoke them. This is an \hl{application of \emph{information hiding}}, a principle that we will encounter again when discussing \emph{distributed objects}.

\highspace
\textcolor{Green3}{\faIcon{question-circle} \textbf{What ``\emph{loosely coupled}'' means in SOA definition?}} A fundamental \textbf{objective} of SOA is to \hl{minimize dependencies between services and their consumers.} For example, if a payment service is used by three different applications (e.g., an e-commerce platform, a mobile app, and a subscription management system), we want to ensure that \hl{internal changes in the payment service do not break the applications that use it.} Therefore, the service's public contract should remain stable, even if its internal implementation evolves over time:
\begin{equation*}
    \text{SOA} = \text{Independent Services} + \text{Well-defined contracts} + \text{Loose Coupling}
\end{equation*}
However, \emph{loose coupling} does not mean \textbf{no coupling at all}. Services or consumers still depend on the service's public contract, availability and behavior. If the service is unavailable or its behavior changes unexpectedly, the consumers may fail to function correctly.

\highspace
In summary, ``\emph{\textbf{loose coupling}}'' in SOA means that \hl{services have \textbf{minimal dependencies on each other},} so one service can change \textbf{independently} without requiring changes to its consumers. It does \textbf{not} mean zero dependencies, as services still rely on each other's public contracts and availability.

\newpage

\begin{flushleft}
    \textcolor{Green3}{\faIcon{sitemap} \textbf{Fundamental Components of SOA}}
\end{flushleft}
A traditional SOA model identifies 4 fundamental concepts.

\begin{figure}[H]
    \centering
    \tikzsetnextfilename{soa-fundamental-components}
    \begin{tikzpicture}[
            font=\normalsize,
            >={Stealth[length=2mm]},
            box/.style={draw=#1, thick, rounded corners=3pt, fill=#1!5},
            head/.style={font=\small\bfseries, text=#1},
            badge/.style={circle, fill=#1, text=white, font=\footnotesize\bfseries,
            inner sep=0pt, minimum size=0.42cm},
            step/.style={circle, draw=#1, fill=white, text=#1, font=\scriptsize\bfseries,
            inner sep=0pt, minimum size=0.36cm},
            flow/.style={->, very thick, #1},
            lbl/.style={font=\footnotesize, align=center, text=#1},
            note/.style={font=\footnotesize\itshape, text=black!65, align=center},
        ]

        % D. Service broker (top)
        \draw[box=orange!70!black] (3.2,5.0) rectangle (8.6,7.0);
        \node[head=orange!70!black] (hb) at (6.1,6.7) {Service broker};
        \node[badge=orange!70!black, left=3pt of hb] {D};
        \node[cylinder, shape border rotate=90, aspect=0.25, draw=orange!70!black, thick,
        fill=white, minimum width=1.3cm, minimum height=0.95cm, font=\footnotesize]
        at (4.25,5.65) {Registry};
        \node[draw=orange!70!black, fill=white, rounded corners=1pt, text width=2.75cm,
            align=left, font=\footnotesize, inner sep=3pt,
        execute at begin node={\hyphenpenalty=10000}] at (6.95,5.65)
        {\textbf{CurrencyService}\\operations, protocol,\\endpoint, interface};
        \node[note, anchor=north] at (5.9,4.92) {discovery only (optional)};

        % C. Service consumer (bottom left)
        \draw[box=blue!50!black] (0,0) rectangle (3.8,3.2);
        \node[head=blue!50!black] (hc) at (2.2,2.9) {Service consumer};
        \node[badge=blue!50!black, left=3pt of hc] {C};
        \node[font=\footnotesize] at (1.9,2.4) {e.g.\ online store};
        \node[draw=black!40, fill=white, rounded corners=1pt, font=\footnotesize\ttfamily,
        align=left, inner sep=3pt] at (1.9,1.6) {convert(100,\\\ \ "USD", "EUR")};
        \node[note] at (1.9,0.45) {knows only the contract};

        % B. Service provider (bottom right) hosting A. the service
        \draw[box=green!40!black] (7.0,0) rectangle (11.4,3.2);
        \node[head=green!40!black] (hp) at (9.45,2.9) {Service provider};
        \node[badge=green!40!black, left=3pt of hp] {B};
        \draw[violet!60!black, thick, dashed, rounded corners=3pt, fill=violet!4]
        (7.2,0.15) rectangle (11.25,2.55);
        \node[badge=violet!60!black] at (7.5,2.25) {A};
        \node[anchor=west, font=\footnotesize\bfseries, text=violet!60!black] at (7.72,2.25)
        {Service};
        \node[anchor=east, font=\footnotesize\itshape, text=violet!60!black] at (11.15,2.25)
        {replicated};
        \foreach \d in {0.16,0.08} {
            \draw[violet!60!black, fill=violet!12, rounded corners=2pt]
            (8.5+\d,0.4+\d) rectangle (10.9+\d,1.6+\d);
        }
        \node[draw=violet!60!black, thick, fill=white, rounded corners=2pt,
        minimum width=2.4cm, minimum height=1.2cm, align=center, font=\footnotesize]
        (svc) at (9.7,1.0) {\textbf{CurrencyService}\\[1pt]implementation\\hidden};
        \node[draw=violet!60!black, thick, fill=violet!25, minimum width=0.42cm,
        minimum height=1.2cm, inner sep=0pt, anchor=east] (api) at (svc.west)
        {\rotatebox{90}{\footnotesize contract}};

        % Interactions: publish - find - bind - invoke
        \draw[flow=green!40!black] (10.6,3.2) |- (8.6,6.1)
        node[pos=0.25, left=4pt, lbl=green!40!black] {publish\\description}
        node[pos=0.25, step=green!40!black] {1};
        \draw[flow=blue!50!black] (0.9,3.2) |- (3.2,6.3)
        node[pos=0.25, left=4pt, lbl=blue!50!black] {find}
        node[pos=0.25, step=blue!50!black] {2};
        \draw[flow=orange!70!black] (3.2,5.4) -| (2.1,3.2)
        node[pos=0.75, right=4pt, lbl=orange!70!black] {endpoint\\+ contract}
        node[pos=0.75, step=orange!70!black] {3};
        \draw[flow=violet!60!black, <->] (3.8,1.0) -- (api.west);
        \node[lbl=violet!60!black, anchor=south] at (5.6,1.05) {bind \& invoke};
        \node[lbl=violet!60!black, anchor=north] at (5.6,0.95) {request / response};
        \node[step=violet!60!black] at (4.15,1.27) {4};
    \end{tikzpicture}
    \caption{\textbf{Fundamental components of SOA.} The provider (B) hosts a service (A) and
        publishes its description to the broker (D) in step~1. The consumer (C) finds the
        service in the registry (2), receives its endpoint and contract (3), and then binds
    to and invokes the service directly at the provider (4).}
    \label{fig:soa-fundamental-components}
\end{figure}

\begin{enumerate}[label={\Alph*.}]
    \item \definition{Service} represents a \textbf{reusable unit of functionality} (e.g., authentication, payment, inventory, shipping services). Ideally, \hl{each service encapsulates its own implementation and exposes a well-defined interface.} A service is a \textbf{logical concept}, not necessarily one process or one physical server. For example, a service can be implemented by multiple servers for scalability and fault tolerance.

    \item \definition{Service Provider} is the component \hl{responsible for \textbf{implementing} and \textbf{exposing} a service.} Its responsibilities include:
        \begin{enumerate}
            \item Implementing the service functionality.
            \item Defining its public interface (contract).
            \item Making the service accessible through a communication endpoint.
            \item Optionally publishing its description to a service broker.
        \end{enumerate}
        \begin{examplebox}[: Currency Conversion Service Provider]
            For example, a company might operate a currency-conversion service accessible through a network API. The provider handles the actual conversion, while consumers only need to know how to request it.
        \end{examplebox}

    \item \definition{Service Consumer} is \hl{any component that \textbf{requires functionality} exposed by another service.}

        \begin{examplebox}[: Currency Conversion Service Consumer]
            For example, an online store might invoke:
\begin{lstlisting}[language=Python]
price_eur = currency_service.convert(
    amount=100,
    source="USD",
    target="EUR"
)
\end{lstlisting}
            This is conceptual pseudocode, as the actual invocation syntax depends on the communication technology used (e.g., REST, gRPC, SOAP). The consumer needs \textbf{to know the service's contract} (i.e., the API), but \textbf{not how} the exchange rate is \textbf{calculated internally}.
        \end{examplebox}

        Importantly, \hl{a component may simultaneously be a provider and a consumer.} For example, an order service might consume a payment service to process payments, while also providing its own order management service to other components. This is similar to the client-server roles, where a server can also act as a client to another server.

    \item\label{def:service-broker} \definition{Service Broker} provides a \hl{mechanism for discovering available services.} It maintains a \hl{registry containing descriptions of services published by providers.}

        \begin{examplebox}[: Currency Conversion Service Discovery]
            Imagine that a consumer needs a service capable of converting currencies. Instead of already knowing the provider's endpoint, it queries the registry to discover a suitable service, and the broker may return information such as the service's identity, its supported operations, its communication protocol, its network endpoint, a description of its interface. Then, the consumer can use this information to communicate with the provider.
        \end{examplebox}

        The broker typically does not execute the service's business logic. Its \hl{main responsibility is \textbf{service discovery}.} Also, a broker is \textbf{not mandatory for every SOA implementation}: services can be configured or discovered through other mechanisms.
\end{enumerate}

\newpage

\begin{flushleft}
    \textcolor{Green3}{\faIcon{exchange-alt} \textbf{Publish-Find-Bind-Invoke Model}}
\end{flushleft}
The interaction among providers, consumers and brokers follows a characteristic sequence of steps, known as the \definition{Publish-Find-Bind-Invoke Model}.

\begin{figure}[H]
    \centering
    \tikzsetnextfilename{soa-publish-find-bind-invoke}
    \begin{tikzpicture}[
            font=\normalsize,
            >={Stealth[length=2mm]},
            actor/.style={draw=#1, thick, fill=#1!12, rounded corners=3pt, minimum width=2.6cm,
            minimum height=0.8cm, align=center, font=\small\bfseries, text=#1},
            badge/.style={circle, fill=#1, text=white, font=\footnotesize\bfseries,
            inner sep=0pt, minimum size=0.42cm},
            life/.style={densely dashed, thick, #1!60},
            act/.style={draw=#1, thick, fill=white},
            phase/.style={fill=#1!7, rounded corners=3pt},
            plabel/.style={rotate=90, font=\footnotesize\bfseries, text=#1},
            call/.style={->, very thick, #1},
            reply/.style={->, very thick, dashed, #1},
            msg/.style={font=\footnotesize, text=#1, align=center},
            step/.style={circle, draw=#1, fill=white, text=#1, font=\scriptsize\bfseries,
            inner sep=0pt, minimum size=0.36cm},
        ]
        \def\xc{2.3}\def\xb{6.2}\def\xp{10.0}

        % phases
        \fill[phase=green!40!black]  (0,6.65) rectangle (11.4,8.05);
        \fill[phase=blue!50!black]   (0,4.15) rectangle (11.4,6.45);
        \fill[phase=violet!60!black] (0,0.45) rectangle (11.4,3.95);
        \node[plabel=green!40!black]  at (0.3,7.35) {Publish};
        \node[plabel=blue!50!black]   at (0.3,5.3)  {Find};
        \node[plabel=violet!60!black] at (0.3,2.2)  {Bind \& invoke};

        % actors and lifelines
        \foreach \x/\c/\n/\l in {\xc/blue!50!black/{Service\\consumer}/C,
            \xb/orange!70!black/{Service\\broker}/D,
        \xp/green!40!black/{Service\\provider}/B} {
            \draw[life=\c] (\x,8.25) -- (\x,0.25);
            \node[actor=\c] (a\l) at (\x,8.7) {\n};
            \node[badge=\c] at (a\l.north west) {\l};
        }

        % activation bars
        \draw[act=orange!70!black] (\xb-0.1,6.95) rectangle (\xb+0.1,7.45);
        \draw[act=orange!70!black] (\xb-0.1,4.65) rectangle (\xb+0.1,5.75);
        \draw[act=blue!50!black]   (\xc-0.1,4.45) rectangle (\xc+0.1,5.95);
        \draw[act=blue!50!black]   (\xc-0.1,0.8) rectangle (\xc+0.1,3.5);
        \draw[act=green!40!black]  (\xp-0.1,1.3) rectangle (\xp+0.1,3.35);

        % 1. publish
        \draw[call=green!40!black] (\xp-0.1,7.2) -- (\xb+0.1,7.2)
        node[midway, above=2pt, msg=green!40!black] {publish service description};
        \node[step=green!40!black] at (\xp+0.45,7.2) {1};

        % 2. find, 3. description + endpoint
        \draw[call=blue!50!black] (\xc+0.1,5.6) -- (\xb-0.1,5.6)
        node[midway, above=2pt, msg=blue!50!black] {find a ``currency\\converter'' service};
        \node[step=blue!50!black] at (\xc-0.45,5.6) {2};
        \draw[reply=orange!70!black] (\xb-0.1,4.8) -- (\xc+0.1,4.8)
        node[midway, below=2pt, msg=orange!70!black] {description + endpoint};
        \node[step=orange!70!black] at (\xb+0.45,4.8) {3};

        % 4. bind and invoke, 5. result
        \draw[call=violet!60!black] (\xc+0.1,3.15) -- (\xp-0.1,3.15)
        node[midway, above=3pt, msg=violet!60!black, fill=violet!7, inner sep=1pt]
        {bind \& invoke: \texttt{convert(100, USD, EUR)}};
        \node[step=violet!60!black] at (\xc-0.45,3.15) {4};
        \node[msg=green!40!black, anchor=east, font=\footnotesize\itshape, align=right]
        at (\xp-0.25,2.3) {executes the\\business logic};
        \node[font=\footnotesize\itshape, text=orange!70!black, fill=violet!7, align=center,
        inner sep=1pt] at (\xb,2.3) {broker not\\involved};
        \draw[reply=green!40!black] (\xp-0.1,1.5) -- (\xc+0.1,1.5)
        node[midway, below=3pt, msg=green!40!black, fill=violet!7, inner sep=1pt]
        {return result: \texttt{92.40 EUR}};
        \node[step=green!40!black] at (\xp+0.45,1.5) {5};
    \end{tikzpicture}
    \caption{\textbf{The publish-find-bind-invoke model.} The provider publishes its service
        description to the broker (1). The consumer finds a suitable service (2) and
        receives its description and endpoint (3). It then binds to the provider and
        invokes the operation directly (4), and the provider returns the result (5);
    the broker is not involved in the invocation.}
    \label{fig:soa-publish-find-bind-invoke}
\end{figure}
\begin{enumerate}[label={\important{Step \arabic*:}}, leftmargin=1.6cm]
    \item \important{Publish.} The \hl{\textbf{provider} registers information describing its service.}
        \begin{examplebox}[: Currency Conversion Service Description]
            For example, a currency conversion service might publish:
\begin{lstlisting}
Service: CurrencyConversion
Operation: convert(amount, source, target)
Endpoint: https://currency.example.com/service
\end{lstlisting}
        \end{examplebox}

    \item \important{Find.} The \hl{\textbf{consumer} requests information about a service} satisfying certain requirements.
        \begin{examplebox}[: Currency Conversion Service Find Request]
            For example, find a service that converts amounts between currencies.
        \end{examplebox}

    \item \important{Discover.} The \hl{\textbf{broker} returns the description of a suitable service,} including the \hl{information} needed \hl{to contact it.}

    \item \important{Bind and Invoke.} The \hl{\textbf{consumer} establishes} the necessary \hl{communication} relationship (\textbf{binding}) and \hl{invokes the operation} through the service's interface.

        \emph{Binding} does not necessarily mean opening a persistent network connection. It means \hl{associating the consumer with the selected service and the protocol details required to invoke it.}

    \item \important{Response.} The \hl{\textbf{provider} executes} the requested operation and \hl{returns the result.}
\end{enumerate}
\textcolor{Green3}{\faIcon{question-circle} \textbf{Why is this useful?}} Suppose the original currency-conversion provider becomes \textbf{unavailable} or is \textbf{replaced} by a different implementation.
\begin{itemize}
    \item[\textcolor{Green3}{\faIcon{check}}] A \textbf{discovery mechanism} can help consumers find an alternative provider exposing a compatible service contract.
    \item[\textcolor{Green3}{\faIcon{check}}] This supports \textbf{flexibility} and \textbf{reduces dependence on specific implementations}.
    \item[\textcolor{Red2}{\faIcon{times}}] However, SOA doesn't automatically provide \textbf{fault tolerance} or \textbf{guarantee} that an \textbf{equivalent provider exists}.
\end{itemize}
Service discovery can contribute to \textbf{location transparency} (\autopageref{item:location-transparency}), a consumer can identify a service logically without embedding the physical location of its implementation. SOA also promotes \textbf{interoperability} (\autopageref{item:interoperability}) and \textbf{modularity} by exposing standardized service contracts.

\newpage
\begin{flushleft}
    \textcolor{Green3}{\faIcon{project-diagram} \textbf{Service Orchestration}}
\end{flushleft}
Until now, we have considered a consumer invoking an individual service. But real distributed applications often need to \hl{combine multiple services to accomplish a single objective.} This introduces the concept of \definition{Service Orchestration}. \textbf{Orchestration} is the process of \hl{coordinating the invocation of multiple services according to a workflow to achieve a particular goal.}

\highspace
\begin{examplebox}[: Processing an Online Order]
    Consider an Order Service that must coordinate three other services to process an online order.

    \begin{center}
        \tikzsetnextfilename{soa-orchestration-online-order}
        \tikzwidth{\textwidth}
        \begin{tikzpicture}[
                font=\normalsize,
                >={Stealth[length=2mm]},
                comp/.style={draw=#1, thick, fill=#1!10, rounded corners=3pt, align=center,
                font=\footnotesize},
                svc/.style={comp=violet!60!black, minimum width=3.0cm, minimum height=1.0cm},
                call/.style={->, very thick, #1},
                reply/.style={->, very thick, dashed, #1},
                step/.style={circle, draw=#1, fill=white, text=#1, font=\scriptsize\bfseries,
                inner sep=0pt, minimum size=0.36cm},
                lbl/.style={font=\footnotesize, text=#1, align=center},
            ]

            % components
            \node[comp=blue!50!black, minimum width=1.5cm, minimum height=1.3cm] (cl)
            at (0.8,2.4) {\small\textbf{Client}};
            \node[comp=green!40!black, minimum width=2.4cm, minimum height=1.8cm] (os)
            at (4.75,2.4) {\small\textbf{OrderService}\\[1pt]\itshape orchestrator:\\
            \itshape runs the workflow};
            \node[svc] (inv) at (9.9,4.25) {\small\textbf{InventoryService}\\check availability,\\reserve product};
            \node[svc] (pay) at (9.9,2.4)  {\small\textbf{PaymentService}\\authorize payment};
            \node[svc] (shp) at (9.9,0.55) {\small\textbf{ShippingService}\\arrange shipment};

            % client <-> orchestrator
            \draw[call=blue!50!black] ([yshift=3mm]cl.east) -- ([yshift=3mm]os.west)
            node[midway, above, lbl=blue!50!black] {place order};
            \draw[reply=green!40!black] ([yshift=-3mm]os.west) -- ([yshift=-3mm]cl.east)
            node[midway, below, lbl=green!40!black] {confirmation};
            \node[step=green!40!black] at ($([yshift=-3mm]cl.east)!0.5!([yshift=-3mm]os.west)$) {4};

            % orchestrator -> services, in workflow order
            \draw[call=violet!60!black] ([yshift=5mm]os.east) -- (inv.west)
            node[pos=0.3, step=violet!60!black] {1};
            \draw[call=violet!60!black] (os.east) -- (pay.west)
            node[pos=0.15, step=violet!60!black] {2}
            node[pos=0.56, below=1pt, lbl=black!60, font=\footnotesize\itshape] {if available};
            \draw[call=violet!60!black] ([yshift=-5mm]os.east) -- (shp.west)
            node[pos=0.2, step=violet!60!black] {3}
            node[pos=0.55, below, sloped, lbl=black!60, font=\footnotesize\itshape] {if authorized};

            % failure handling
            \node[comp=red!60!black, minimum width=2.7cm, font=\footnotesize] (fail) at (4.75,0.3)
            {payment fails $\Rightarrow$ cancel\\order, release reservation};
            \draw[reply=red!60!black] (os.south) -- (fail.north);
        \end{tikzpicture}
    \end{center}

    When a customer places an order, the orchestrator \texttt{OrderService} might execute the following workflow:
    \begin{enumerate}
        \item Invoke \texttt{InventoryService} to check product availability.
        \item If the product is available, invoke \texttt{PaymentService} to authorize the payment.
        \item If authorization succeeds, reserve the product and arrange shipment through the appropriate services (e.g., \texttt{ShippingService}).
        \item Return the order confirmation.
    \end{enumerate}
    If the payment fails, the orchestrator must decide what happens next, such as cancelling the order or releasing a reservation.

    \highspace
    The important point is that \hl{each service implements a specific functionality, while the orchestrator \textbf{coordinates} how those functionalities contribute to a larger operation.}
\end{examplebox}

\newpage

\begin{flushleft}
    \textcolor{Green3}{\faIcon{globe} \textbf{Traditional Web Services}}
\end{flushleft}
We have discussed SOA as an architectural style, independently of the technology used to implement it. Now we examine \textbf{one of its historically important implementations}: \definition{SOAP-based Web Services}.

\highspace
In general, a \definition{Web Service} is a \hl{software system that enables interoperable machine-to-machine communication over a network.} Traditional Web Services typically \textbf{rely on three standards}:
\begin{itemize}
    \item \definition{Web Services Description Language (WSDL)} \hl{specifies the \textbf{service contract}}, like the operations it exposes, the messages it can send and receive, and the communication details (e.g., protocols, endpoints).

        Specifically, WSDL is an \hl{XML-based language used to \textbf{describe the public interface} of a Web Service.} It defines information such as:
        \begin{itemize}
            \item Which operations the service provides.
            \item Which messages and data types are accepted.
            \item Which messages are returned.
            \item Which protocol and message format are used.
            \item Where the service can be accessed.
        \end{itemize}
        So in general, \hl{WSDL tells the consumer \textbf{what the service exposes} and \textbf{how to invoke it}.}

        \begin{examplebox}[: Calculator Service Contract]
            Consider a simple calculator Web Service. Its interface might conceptually define:
\begin{lstlisting}
CalculatorService

Operation:
    add(a: integer, b: integer) -> integer

Input:
    a, b

Output:
    result
\end{lstlisting}
            A \textbf{WSDL document} describes this contract in a \textbf{standardized}, \textbf{machine-readable} way. The standardized description allows different systems, like a Java client and a C\# service, to communicate without needing to know each other's internal implementations. The consumer only needs to understand the interface and communication conventions. Furthermore, development tools can even generate client-side proxy code from a WSDL description.
        \end{examplebox}

        \newpage

    \item \definition{Simple Object Access Protocol (SOAP)} defines an \textbf{XML-based messaging protocol} for exchanging requests and responses.

        In general, WSDL describes the interface, but \textbf{we still need to exchange actual invocation messages}. This is the role of \textbf{SOAP}, which is an \hl{XML-based messaging protocol defining \textbf{how structured messages are exchanged between communicating components}.}

        \begin{examplebox}[: SOAP Request and Response]
            For example, a \textbf{consumer} might invoke the calculator's \texttt{add} operation by \textbf{sending an illustrative SOAP request}:
\begin{lstlisting}[language=xml]
<soap:Envelope
    xmlns:soap="http://www.w3.org/2003/05/soap-envelope">

    <soap:Body>
        <add xmlns="http://example.com/calculator">
            <a>5</a>
            <b>3</b>
        </add>
    </soap:Body>

</soap:Envelope>
\end{lstlisting}
            The \textbf{service executes} the operation and \textbf{returns a corresponding SOAP response}:
\begin{lstlisting}[language=xml]
<soap:Envelope
    xmlns:soap="http://www.w3.org/2003/05/soap-envelope">

    <soap:Body>
        <addResponse xmlns="http://example.com/calculator">
            <result>8</result>
        </addResponse>
    </soap:Body>

</soap:Envelope>
\end{lstlisting}
            Notice the structure:
            \begin{itemize}
                \item \texttt{Envelope} identifies the message as a SOAP message.
                \item \texttt{Header} (optional) carries additional information, potentially including security or routing metadata.
                \item \texttt{Body} contains the application-specific message content.
            \end{itemize}
            Note that SOAP specifies the \textbf{messaging structure} and \textbf{processing rules}, but it does not determine the application's business logic (i.e., what the service actually does).
        \end{examplebox}

    \item \definition{Universal Description, Discovery, and Integration (UDDI)} defines a \textbf{registry framework} for publishing and discovering services.

        Precisely, UDDI is a \hl{standardized framework for \textbf{publishing} and \textbf{discovering Web Services} through registries.} It implements the \emph{service broker} concept introduced earlier on \autopageref{def:service-broker}.

        In other words, UDDI tells consumers \hl{which services are available and where to find information about them.} UDDI is historically important for understanding SOA, but many real-world Web Services are accessed through known endpoints or other discovery mechanisms without relying on UDDI registries.
\end{itemize}

\highspace
\begin{table}[!htp]
    \centering
    \begin{tabular}{@{} l c l @{}}
        \toprule
        \textbf{SOA activity} & \textbf{Technology} & \textbf{Responsibility} \\
        \midrule
        Describe & WSDL & Defines the service contract \\[.5em]
        Publish / Discover & UDDI & Registers and discovers services \\[.5em]
        Bind / Invoke & SOAP & Exchanges invocation messages \\
        \bottomrule
    \end{tabular}
    \caption{\textbf{Traditional Web Services technologies.}}
\end{table}

\highspace
\begin{figure}[H]
    \centering
    \tikzsetnextfilename{traditional-web-services}
    \begin{tikzpicture}[
            font=\normalsize,
            >={Stealth[length=2mm]},
            panel/.style={draw=#1, thick, rounded corners=3pt, fill=#1!4},
            headbar/.style={fill=#1!18, rounded corners=2pt},
            acr/.style={font=\large\bfseries, text=#1},
            verb/.style={draw=#1, fill=white, rounded corners=6pt, font=\footnotesize\bfseries,
            text=#1, inner xsep=4pt, inner ysep=2pt, text height=1.6ex, text depth=0.3ex},
            expand/.style={font=\footnotesize\itshape, text=black!75, align=center},
            idea/.style={font=\footnotesize, align=center, text width=3.3cm,
            execute at begin node={\hyphenpenalty=10000}},
            code/.style={font=\footnotesize\ttfamily, align=left},
            part/.style={font=\footnotesize\bfseries, text=teal!60!black, anchor=west},
            val/.style={font=\footnotesize\ttfamily, anchor=west},
        ]
        % ---------- panels and headers
        \foreach \x/\c/\a/\v/\e in {0/teal!60!black/WSDL/Describe/{Web Services\\Description Language},
            3.875/violet!60!black/SOAP/Communicate/{Simple Object\\Access Protocol},
        7.75/orange!70!black/UDDI/Discover/{Universal Description,\\Discovery and Integration}} {
            \draw[panel=\c] (\x,0) rectangle (\x+3.65,7.45);
            \fill[headbar=\c] (\x+0.08,5.85) rectangle (\x+3.57,7.37);
            \node[acr=\c] at (\x+1.825,6.98) {\a};
            \node[expand] at (\x+1.825,6.28) {\e};
            \draw[\c, thick] (\x+0.5,6.67) -- (\x+3.15,6.67);
            \node[verb=\c] at (\x+1.825,5.4) {\v};
        }

        % ---------- WSDL: the contract document, with its four parts
        \draw[teal!60!black, thick, fill=white]
        (0.3,1.08) -- (3.35,1.08) -- (3.35,4.3) -- (3.05,4.6) -- (0.3,4.6) -- cycle;
        \draw[teal!60!black, thick] (3.05,4.6) -- (3.05,4.3) -- (3.35,4.3);
        \fill[teal!60!black!15] (0.33,4.0) rectangle (3.02,4.57);
        \node[font=\footnotesize\ttfamily\bfseries, text=teal!60!black] at (1.67,4.29)
        {calculator.wsdl};
        \foreach \y/\p/\v in {3.74/types/{a, b, result:\,int},
            3.03/interface/{add(a, b)},
            2.32/binding/{SOAP over HTTP},
        1.61/service/{endpoint URL}} {
            \node[part] at (0.4,\y) {\p};
            \node[val] at (0.55,\y-0.31) {\v};
        }
        \foreach \y in {3.25,2.54,1.83}
        \draw[teal!60!black!30] (0.45,\y) -- (3.2,\y);

        % ---------- SOAP: envelope carried by HTTP
        \draw[black!45, thick, dashed, rounded corners=2pt] (4.05,1.15) rectangle (7.35,4.75);
        \node[anchor=north west, font=\footnotesize\itshape, text=black!60] at (4.1,4.75)
        {HTTP transport};
        \draw[violet!60!black, thick, fill=white, rounded corners=2pt] (4.2,1.3) rectangle (7.2,4.35);
        \node[anchor=north west, font=\footnotesize\bfseries, text=violet!60!black]
        at (4.25,4.33) {Envelope};
        \draw[violet!60!black, dashed, fill=violet!6] (4.35,3.35) rectangle (7.05,3.85);
        \node[font=\footnotesize, text=violet!60!black] at (5.7,3.6) {Header (optional)};
        \draw[violet!60!black, fill=violet!12] (4.35,1.45) rectangle (7.05,3.2);
        \node[anchor=north west, font=\footnotesize\bfseries, text=violet!60!black]
        at (4.4,3.18) {Body};
        \node[code, anchor=north west] at (4.75,2.78)
        {<add>\\[-1pt]\ \ <a>5</a>\\[-1pt]\ \ <b>3</b>\\[-1pt]</add>};

        % ---------- UDDI: the registry (service broker)
        \node[cylinder, shape border rotate=90, aspect=0.25, draw=orange!70!black, thick,
            fill=white, minimum width=2.6cm, minimum height=1.3cm, align=center,
        font=\footnotesize] (reg) at (9.575,3.85) {\textbf{Registry}\\(service broker)};
        \node[draw=red!60!black, fill=red!8, rounded corners=2pt, font=\footnotesize,
        minimum width=1.5cm, text height=1.6ex, text depth=0.4ex] (prov) at (8.65,1.6) {provider};
        \node[draw=blue!50!black, fill=blue!10, rounded corners=2pt, font=\footnotesize,
        minimum width=1.5cm, text height=1.6ex, text depth=0.4ex] (cons) at (10.5,1.6) {consumer};
        \draw[->, thick, red!60!black] (prov.north) -- (prov.north |- reg.south)
        node[midway, fill=orange!4, inner sep=1pt, font=\footnotesize] {publish};
        \draw[->, thick, blue!50!black] (cons.north) -- (cons.north |- reg.south)
        node[midway, fill=orange!4, inner sep=1pt, font=\footnotesize] {find};

        % ---------- key ideas
        \node[idea] at (1.825,0.62) {\emph{what} the service offers and \emph{how} to call it};
        \node[idea] at (5.7,0.62)   {XML messages carrying requests and responses};
        \node[idea] at (9.575,0.62) {\emph{which} services exist and \emph{where} to find them};
    \end{tikzpicture}
    \caption{\textbf{The three standards of traditional Web Services.} WSDL describes the
        service contract (types, interface, binding, endpoint), SOAP carries request and response messages (usually over
        HTTP), and UDDI provides a registry where providers publish and consumers find
    services.}
    \label{fig:traditional-web-services}
\end{figure}

\newpage

\begin{examplebox}[: How WSDL, SOAP, and UDDI Work Together]
    Consider a Java application that needs a remote calculator.
    \begin{enumerate}
        \item The calculator provider makes its service description available through WSDL.
        \item The provider may register its service in a UDDI registry.
        \item The Java application discovers the calculator service through the registry.
        \item It retrieves the WSDL description and configures or generates a client.
        \item The client sends a SOAP request to invoke \texttt{add(5, 3)}.
        \item The provider processes the request and responds with the result \texttt{8}.
    \end{enumerate}
    Note that \hl{the broker is involved in discovery, not necessarily in every operation invocation.}

    \begin{center}
        \tikzsetnextfilename{wsdl-soap-uddi-together}
        \tikzwidth{\textwidth}
        \begin{tikzpicture}[
                font=\normalsize,
                >={Stealth[length=2mm]},
                box/.style={draw=#1, thick, rounded corners=3pt, fill=#1!5},
                head/.style={font=\small\bfseries, text=#1, align=center},
                inner/.style={draw=#1, thick, fill=white, rounded corners=2pt, align=center,
                font=\footnotesize},
                flow/.style={->, very thick, #1},
                reply/.style={->, very thick, dashed, #1},
                lbl/.style={font=\footnotesize, align=center, text=#1},
                step/.style={circle, solid, draw=#1, fill=white, text=#1, font=\scriptsize\bfseries,
                inner sep=0pt, minimum size=0.36cm},
            ]

            % UDDI registry (broker)
            \draw[box=orange!70!black] (3.3,4.55) rectangle (7.7,6.65);
            \node[head=orange!70!black] at (5.5,6.32) {UDDI registry};
            \node[cylinder, shape border rotate=90, aspect=0.25, draw=orange!70!black, thick,
            fill=white, minimum width=1.3cm, minimum height=0.9cm, font=\footnotesize]
            at (4.2,5.22) {Registry};
            \node[inner=orange!70!black, align=left, inner sep=3pt] at (6.35,5.22)
            {\textbf{Calculator}\\$\rightarrow$ WSDL location};

            % WSDL document
            \draw[teal!60!black, thick, fill=white]
            (9.15,4.7) -- (10.75,4.7) -- (10.75,5.95) -- (10.5,6.2) -- (9.15,6.2) -- cycle;
            \draw[teal!60!black, thick] (10.5,6.2) -- (10.5,5.95) -- (10.75,5.95);
            \node[font=\footnotesize, text=teal!60!black, align=center] (wsdl) at (9.95,5.45)
            {\textbf{WSDL}\\\ttfamily add(a,b)};

            % Java application (consumer)
            \draw[box=blue!50!black] (0,0) rectangle (3.7,3.1);
            \node[head=blue!50!black] at (1.85,2.68) {Java application\\[-2pt]
            {\footnotesize\mdseries\itshape service consumer}};
            \node[inner=violet!60!black, minimum width=2.8cm, minimum height=0.85cm] (proxy)
            at (1.85,1.45) {\textbf{client proxy}\\generated from WSDL};

            % Calculator provider
            \draw[box=red!60!black] (7.7,0) rectangle (11.4,3.1);
            \node[head=red!60!black] at (9.55,2.68) {Calculator provider\\[-2pt]
            {\footnotesize\mdseries\itshape service provider}};
            \node[inner=red!60!black, minimum width=2.8cm, minimum height=0.85cm] (svc)
            at (9.55,1.45) {\textbf{CalculatorService}\\\texttt{add(a, b)}};

            % 1. describe with WSDL, 2. register in UDDI
            \draw[flow=teal!60!black] (9.95,3.1) -- (9.95,4.7)
            node[pos=0.5, step=teal!60!black] {1}
            node[pos=0.5, left=6pt, lbl=teal!60!black] {describe};
            \draw[flow=orange!70!black] (9.15,5.75) -- (7.7,5.75)
            node[pos=0.5, step=orange!70!black] {2}
            node[pos=0.5, above=5pt, lbl=orange!70!black] {register};

            % 3. discover, 4. retrieve WSDL
            \draw[flow=blue!50!black] (0.8,3.1) |- (3.3,5.75)
            node[pos=0.25, step=blue!50!black] {3}
            node[pos=0.25, right=6pt, lbl=blue!50!black] {discover};
            \draw[flow=teal!60!black] (9.15,4.85) -- (2.7,3.1)
            node[pos=0.13, step=teal!60!black] {4}
            node[pos=0.55, below, sloped, lbl=teal!60!black] {retrieve WSDL, generate client};

            % 5. SOAP request, 6. SOAP response (direct, no registry)
            \draw[flow=violet!60!black] (proxy.east |- 0,1.7) -- (svc.west |- 0,1.7)
            node[pos=0.12, step=violet!60!black] {5}
            node[pos=0.55, above, lbl=violet!60!black] {SOAP request \texttt{add(5, 3)}};
            \draw[reply=red!60!black] (svc.west |- 0,1.2) -- (proxy.east |- 0,1.2)
            node[pos=0.88, step=red!60!black] {6}
            node[pos=0.45, below, lbl=red!60!black] {SOAP response \texttt{8}};
            \node[font=\footnotesize\itshape, text=black!60, align=center] at (5.7,0.4)
            {direct call,\\UDDI not involved};
        \end{tikzpicture}
        \captionof{figure}{\textbf{How WSDL, SOAP and UDDI work together.} The provider describes its
            calculator service with WSDL (1) and registers it in UDDI (2). The Java
            application discovers it (3), retrieves the WSDL and generates a client (4), then
            invokes \texttt{add(5, 3)} with a SOAP request (5) and receives \texttt{8} in the
            SOAP response (6).
        The SOAP exchange goes directly from consumer to provider, without UDDI.}
        \label{fig:wsdl-soap-uddi-together}
    \end{center}
\end{examplebox}

\newpage
\begin{flushleft}
    \textcolor{Green3}{\faIcon{balance-scale} \textbf{SOA vs. Client-Server and Multi-Tier Architectures}}
\end{flushleft}
After introducing SOA and traditional Web Services, a natural question arises: \emph{\textbf{isn't SOA simply another name for client-server?}} This question is understandable, as both SOA and client-server architectures involve a consumer (client) invoking a service (server). However, there are important distinctions because \textbf{emphasize different aspects of a system}, even though they can coexist.

\highspace
\begin{table}[H]
    \centering
    \begin{tabular}{@{} l p{20em} @{}}
        \toprule
        \textbf{Architecture} & \textbf{Main question} \\
        \midrule
        Client-Server & \emph{Who requests functionality, and who provides it?} \\[1.4em]
        Multi-Tier & \emph{How are application responsibilities distributed across execution tiers?} \\[1.4em]
        SOA & \emph{How can a system be composed of loosely coupled, reusable services?} \\
        \bottomrule
    \end{tabular}
\end{table}

\highspace
For example, consider the classic e-commerce example:
\begin{itemize}
    \item It uses \textbf{client-server} interactions because consumers (i.e., customers) invoke providers (i.e., servers) to request and receive services.
    \item It can be \textbf{three-tier} because presentation (i.e., web browser), application logic (i.e., order processing), and data management (i.e., database) are separated into different tiers.
    \item It follows \textbf{SOA principles} because authentication, payment, inventory, and shipping are exposed as independently defined services.
    \item It might use \textbf{SOAP-based Web Services} as the communication technology between some components (e.g., between the order processing service and the payment service).
\end{itemize}
So, SOA is not a replacement for client-server or multi-tier architectures. Also, \hl{SOA is not synonymous with SOAP}, because SOAP-based Web Services are \textbf{just one way to implement SOA}. The broader SOA principles can also be implemented with other communication technologies.

\highspace
\begin{takeawaysbox}
    \begin{enumerate}
        \item \textbf{SOA organizes distributed applications around services.} Services are reusable, loosely coupled units of functionality exposed through well-defined interfaces.

        \item \textbf{Four fundamental concepts define the traditional model.}
            \begin{itemize}
                \item \textbf{Service}: the functionality being offered.
                \item \textbf{Provider}: implements and exposes the service.
                \item \textbf{Consumer}: discovers and invokes the service.
                \item \textbf{Broker}: maintains information used for service discovery.
            \end{itemize}

        \item \textbf{The typical SOA interaction follows a publish-find-bind-invoke model.}

        \item \textbf{Orchestration combines services into workflows.} An orchestrator coordinates multiple service invocations to accomplish a larger objective.

        \item \textbf{Remember the three traditional Web Services technologies.}
            \begin{equation*}
                \begin{aligned}
                    \text{WSDL} &\rightarrow \text{Describe}\\
                    \text{SOAP} &\rightarrow \text{Communicate}\\
                    \text{UDDI} &\rightarrow \text{Discover}
                \end{aligned}
            \end{equation*}
    \end{enumerate}
    In general, \hl{SOA changes the way we think about building distributed applications.} Instead of thinking only about distributing application components across machines, we think about \hl{composing independently defined services into larger applications.} This makes the service contract, loose coupling, discovery, and composition central architectural concerns.
\end{takeawaysbox}